%----------------------------------------------------------------------------------------
% LIMIETEN EN CONTINUÏTEIT. Alle opgaven zijn eigen werk.
%
% Opbouw: Start, leerstof (theorie, werkwijze of uitleg, voorbeeld, korte oefening), en op een
% nieuwe pagina de Opdrachten (door elkaar, stijgende moeilijkheid), Koppel en Test jezelf.
% Labels van dit hoofdstuk beginnen met h1-.
%----------------------------------------------------------------------------------------

\mychapter{1}{Limieten en continuïteit}

\lkr{Leerplan B+S: LPD 18-21 (afgeleiden), het deel over het limietbegrip, rekenregels, onbepaalde vormen, asymptoten als limieten en continuïteit. De $\varepsilon$-$\delta$-definitie staat erin, met één uitgewerkt bewijs voor een eerstegraadsfunctie; in dit pakket kan ze wellicht weg (na te kijken met de leerplantekst), en dan ook de bewijzen van de rekenregels, die erop steunen. In Rationale functies werden asymptoten informeel ingevoerd: hier krijgen ze een definitie.}

Wat gebeurt er met $\frac{x^2 - 1}{x - 1}$ als $x$ naar $1$ gaat? Invullen kan niet: dat geeft $\frac00$. Toch komen de functiewaarden steeds dichter bij $2$. Een limiet beschrijft precies dat: waar de functiewaarden naartoe gaan, ook als de functie er zelf niet bestaat. Limieten zijn de basis van de afgeleide, die in hoofdstuk 2 volgt.

\begin{oefeningen}[Start]
\oef{h1-start-1}

\oef{h1-verken}
\end{oefeningen}

%========================================================================================
\section{Het begrip limiet}

\subsection{Waar gaan de functiewaarden naartoe?}

\naast{%
Bekijk $f(x) = \frac{x^2 - 1}{x - 1}$ voor $x$ dicht bij $1$:

\vspace{1mm}
{\renewcommand{\arraystretch}{1.2}%
\begin{tabular}{c|cccc}
$x$ & $0{,}9$ & $0{,}99$ & $1{,}01$ & $1{,}1$ \\ \hline
$f(x)$ & $1{,}9$ & $1{,}99$ & $2{,}01$ & $2{,}1$
\end{tabular}}

\vspace{1mm}
Voor $x \neq 1$ is $f(x) = \frac{(x - 1)(x + 1)}{x - 1} = x + 1$: de grafiek is een rechte met een gat in $(1, 2)$. Hoe dichter $x$ bij $1$ komt, langs links of langs rechts, hoe dichter $f(x)$ bij $2$ komt. We schrijven $\lim_{x \to 1} f(x) = 2$.}{%
\begin{tikzpicture}[x=.6cm,y=.6cm,font=\scriptsize]
  \rooster{-1}{1}{3}{-1}{1}{4}
  \draw[accent,thick] (-1,0) -- (3,4);
  \open{(1,2)}
  \draw[grijs,dashed] (1,0) -- (1,2) (0,2) -- (1,2);
\end{tikzpicture}}

Een limiet zegt dus niets over de functiewaarde in $1$ zelf, maar alles over de functiewaarden \textit{rondom} $1$.

\begin{theorie}
$\formule{\lim_{x \to a} f(x) = L}$ betekent: $f(x)$ komt willekeurig dicht bij $L$ als $x$ voldoende dicht bij $a$ komt, met $x \neq a$. We zeggen: de \textbf{\textcolor{accent}{limiet}} van $f(x)$ voor $x$ naar $a$ is $L$.
\end{theorie}

De waarde $f(a)$ speelt geen rol: $x$ komt willekeurig dicht bij $a$, maar is nooit gelijk aan $a$. De functie hoeft in $a$ zelfs niet te bestaan, zoals hierboven. En bestaat $f(a)$ wel, dan kan die waarde iets anders zijn dan de limiet.

\voorbeeld{}{De figuur toont de grafiek van een functie $f$. Lees $\lim_{x \to 2} f(x)$ en $f(2)$ af.\label{vb:h1-aflezen}}{%
\naast{%
\uitleg{Kijk naar de grafiek links en rechts van $x = 2$, en apart naar het punt boven $2$ zelf.}
\begin{stappenvrij}{3.4cm}{6.2cm}
\lim_{x \to 2} f(x) = 3 & Schuif langs de grafiek naar $x = 2$, van links en van rechts: de hoogte nadert $3$, de hoogte van het open bolletje. \\
f(2) = 1 & Het volle bolletje zegt welke waarde $f$ in $2$ echt aanneemt.
\end{stappenvrij}
De limiet en de functiewaarde verschillen hier. Dat is geen tegenspraak: de limiet kijkt alleen naar $x \neq 2$.}{%
\begin{tikzpicture}[x=.55cm,y=.55cm,font=\scriptsize]
  \rooster{-1}{1}{4}{-1}{1}{4}
  \draw[accent,thick] (-1,0) -- (3.5,4.5);
  \open{(2,3)} \vol{(2,1)}
\end{tikzpicture}}}

\begin{oefeningen}
\oef{h1-aflezen}
\end{oefeningen}

\subsection{De definitie met \texorpdfstring{$\varepsilon$}{epsilon} en \texorpdfstring{$\delta$}{delta}}

\naast{%
Wat betekent \textit{willekeurig dicht} precies? Bekijk het als een uitdaging. Iemand kiest een nauwkeurigheid $\varepsilon > 0$: de functiewaarden moeten in de strook tussen $L - \varepsilon$ en $L + \varepsilon$ liggen. Jij moet dan een afstand $\delta > 0$ vinden zodat dat lukt voor alle $x \neq a$ tussen $a - \delta$ en $a + \delta$.

Lukt dat voor elke $\varepsilon$, hoe klein ook, dan is $L$ de limiet. Een kleinere $\varepsilon$ vraagt meestal een kleinere $\delta$.}{%
\begin{tikzpicture}[x=1.05cm,y=1.05cm,font=\scriptsize]
  \fill[accent!12] (0,1.5) rectangle (3.9,2.5);
  \fill[black!8] (1.5,0) rectangle (2.5,3.7);
  \draw[->,grijs] (-.2,0) -- (4.1,0) node[right,text=black] {$x$};
  \draw[->,grijs] (0,-.2) -- (0,3.9) node[above,text=black] {$y$};
  \draw[accent,dashed] (0,2.5) -- (3.9,2.5) (0,1.5) -- (3.9,1.5);
  \draw[grijs,dashed] (1.5,0) -- (1.5,3.7) (2.5,0) -- (2.5,3.7);
  \kromme{.3}{3.6}{2+.8*(\x-2)+.15*(\x-2)^2}
  \open{(2,2)}
  \node[left] at (0,2.5) {$L + \varepsilon$};
  \node[left] at (0,2) {$L$};
  \node[left] at (0,1.5) {$L - \varepsilon$};
  \node[below left,inner xsep=0pt] at (1.5,0) {$a - \delta$};
  \node[below] at (2,0) {$a$};
  \node[below right,inner xsep=0pt] at (2.5,0) {$a + \delta$};
\end{tikzpicture}}

\begin{theorie}
$\formule{\lim_{x \to a} f(x) = L}$ betekent nauwkeurig: voor elke $\varepsilon > 0$ bestaat er een $\delta > 0$ zodat
\[\formule{0 < |x - a| < \delta \Rightarrow |f(x) - L| < \varepsilon}.\]
\end{theorie}

Lees de formule als de figuur: $|x - a| < \delta$ is de grijze strook, $|f(x) - L| < \varepsilon$ de gekleurde. De voorwaarde $0 < |x - a|$ zegt dat $x \neq a$: de waarde in $a$ zelf telt niet mee.

Met de definitie bewijs je dat een limiet klopt. Je toont dan aan dat er bij elke $\varepsilon$ een geschikte $\delta$ bestaat, en dat doe je door $\delta$ uit te drukken in $\varepsilon$.

\voorbeeld{}{Bewijs met de definitie: $\lim_{x \to 2} (3x - 1) = 5$.\label{vb:h1-definitie}}{%
\uitleg{We kunnen geen getal kiezen voor $\varepsilon$: het bewijs moet voor elke $\varepsilon$ werken. Neem daarom een willekeurige $\varepsilon > 0$. Schrijf $|f(x) - L|$ uit tot je de afstand $|x - 2|$ ziet verschijnen.}
\begin{stappenbreed}
|(3x - 1) - 5| = |3x - 6| = 3|x - 2| & We werken de afstand tussen $f(x)$ en $L$ uit. Er verschijnt $3$ keer de afstand tussen $x$ en $2$. \\[1ex]
3|x - 2| < \varepsilon \iff |x - 2| < \tfrac{\varepsilon}{3} & We willen dat dit kleiner is dan $\varepsilon$. Dat lukt zodra $x$ dichter dan $\frac{\varepsilon}{3}$ bij $2$ ligt. \\[1ex]
\delta = \tfrac{\varepsilon}{3} & Dat wordt onze $\delta$. Hij hangt af van $\varepsilon$, zoals het hoort. \\[1ex]
0 < |x - 2| < \delta \Rightarrow |f(x) - 5| = 3|x - 2| < 3\delta = \varepsilon & Controle: met deze $\delta$ geldt de voorwaarde uit de definitie. Omdat $\varepsilon$ willekeurig was, is de limiet bewezen.
\end{stappenbreed}
Met getallen: wie $\varepsilon = 0{,}03$ eist, krijgt $\delta = 0{,}01$. Voor $1{,}99 < x < 2{,}01$ ligt $3x - 1$ inderdaad tussen $4{,}97$ en $5{,}03$.}

\begin{werkwijze}
Neem een willekeurige $\varepsilon > 0$. Werk $|f(x) - L|$ uit tot er een veelvoud $c\,|x - a|$ staat. Kies $\delta = \frac{\varepsilon}{c}$. Controleer: uit $0 < |x - a| < \delta$ volgt $|f(x) - L| < c\,\delta = \varepsilon$.
\end{werkwijze}

\begin{oefeningen}
\oef{h1-begrip}
\end{oefeningen}

\subsection{Linker- en rechterlimiet}

Soms komt $f(x)$ langs links bij een ander getal uit dan langs rechts. Dan bekijk je de twee kanten apart.

\begin{theorie}
De \textbf{\textcolor{accent}{linkerlimiet}} $\formule{\liml{x}{a} f(x)}$ gebruikt alleen $x < a$, de \textbf{\textcolor{accent}{rechterlimiet}} $\formule{\limr{x}{a} f(x)}$ alleen $x > a$. In de definitie met $\varepsilon$ en $\delta$ vervang je $0 < |x - a| < \delta$ door $\formule{a - \delta < x < a}$ voor de linkerlimiet, en door $\formule{a < x < a + \delta}$ voor de rechterlimiet.

\begin{stelling}{Limiet en eenzijdige limieten}
$\formule{\lim_{x \to a} f(x) = L \iff \liml{x}{a} f(x) = L \text{ en } \limr{x}{a} f(x) = L}$
\end{stelling}
\begin{bewijs}
$\Rightarrow$\ Geldt $|f(x) - L| < \varepsilon$ voor alle $x$ met $0 < |x - a| < \delta$, dan zeker voor de $x$ links van $a$ en voor de $x$ rechts van $a$.

$\Leftarrow$\ Neem $\varepsilon > 0$. De linkerlimiet geeft een $\delta_1$ en de rechterlimiet een $\delta_2$ met
\bewijsregel{a - \delta_1 < x < a \Rightarrow |f(x) - L| < \varepsilon \qquad\text{en}\qquad a < x < a + \delta_2 \Rightarrow |f(x) - L| < \varepsilon .}
Neem $\delta$ de kleinste van $\delta_1$ en $\delta_2$. Is $0 < |x - a| < \delta$, dan ligt $x$ links van $a$ binnen $\delta_1$ of rechts van $a$ binnen $\delta_2$. In beide gevallen is $|f(x) - L| < \varepsilon$.
\end{bewijs}
\end{theorie}

De limiet bestaat dus alleen als beide kanten naar hetzelfde getal gaan. Gaan ze naar verschillende getallen, of bestaat een van beide niet, dan is er geen limiet.

\voorbeeld{}{Bereken $\liml{x}{2} f(x)$ en $\limr{x}{2} f(x)$ voor $f(x) = \dfrac{|x - 2|}{x - 2}$. Bestaat $\lim_{x \to 2} f(x)$?\label{vb:h1-sprong}}{%
\naast{%
\uitleg{De absolute waarde hangt af van het teken van $x - 2$. Daarom splitsen we in links en rechts van $2$.}
\begin{stappenvrij}{3.7cm}{5.9cm}
{f(x) = \frac{-(x - 2)}{x - 2} = -1} & Voor $x < 2$. Links van $2$ is $x - 2 < 0$, dus $|x - 2| = -(x - 2)$. \\
\liml{x}{2} f(x) = -1 & De functie is daar constant $-1$. \\
{f(x) = \frac{x - 2}{x - 2} = 1} & Voor $x > 2$: rechts is $x - 2 > 0$. \\
\limr{x}{2} f(x) = 1 & \\
{\lim_{x \to 2} f(x) \text{ bestaat niet}} & Links en rechts verschillen.
\end{stappenvrij}}{%
\begin{tikzpicture}[x=.55cm,y=.55cm,font=\scriptsize]
  \rooster{-1}{1}{4}{-2}{1}{2}
  \draw[accent,thick] (-1,-1) -- (2,-1);
  \draw[accent,thick] (2,1) -- (4.3,1);
  \open{(2,-1)} \open{(2,1)}
\end{tikzpicture}}}

\begin{oefeningen}
\oef{h1-links-rechts}
\end{oefeningen}

%========================================================================================
\section{Limieten berekenen}

Elke limiet met $\varepsilon$ en $\delta$ bewijzen is veel werk. Gelukkig hoeft dat niet: een paar rekenregels, één keer bewezen, volstaan voor de meeste limieten. We bouwen ze stap voor stap op, vanaf de eenvoudigste functies.

\subsection{Twee basislimieten}

De eenvoudigste functies zijn een constante en de functie $x$ zelf. Hun limiet is wat je verwacht.

\begin{theorie}
\begin{stelling}{Basislimieten}
Voor elk getal $c$ en elk getal $a$ is \quad $\formule{\lim_{x \to a} c = c}$ \quad en \quad $\formule{\lim_{x \to a} x = a}$.
\end{stelling}
\begin{bewijs}
Neem $\varepsilon > 0$. Voor de constante is $|c - c| = 0 < \varepsilon$ voor elke $x$, dus elke $\delta$ werkt. Voor de functie $x$ is $|f(x) - a| = |x - a|$; neem $\delta = \varepsilon$, dan
\bewijsregel{0 < |x - a| < \delta \Rightarrow |x - a| < \varepsilon .}
\end{bewijs}
\end{theorie}

Met deze twee limieten en de rekenregels hieronder vind je de limiet van elke veelterm.

\subsection{Limiet van een som}

Komen $f(x)$ en $g(x)$ dicht bij $L$ en $M$, dan komt hun som dicht bij $L + M$. Dat is de eerste rekenregel.

\begin{theorie}
\begin{stelling}{Somregel}
Als $\lim_{x \to a} f(x) = L$ en $\lim_{x \to a} g(x) = M$, dan is
\[\formule{\lim_{x \to a} \bigl(f(x) + g(x)\bigr) = L + M} \qquad\text{en}\qquad \formule{\lim_{x \to a} \bigl(f(x) - g(x)\bigr) = L - M}.\]
\end{stelling}
\begin{bewijs}
Neem $\varepsilon > 0$. Pas de definitie toe met $\frac{\varepsilon}{2}$: er is een $\delta_1$ zodat $|f(x) - L| < \frac{\varepsilon}{2}$ voor $0 < |x - a| < \delta_1$, en een $\delta_2$ zodat $|g(x) - M| < \frac{\varepsilon}{2}$ voor $0 < |x - a| < \delta_2$. Neem $\delta$ de kleinste van de twee. Met de driehoeksongelijkheid $|u + v| \leqslant |u| + |v|$ is dan voor $0 < |x - a| < \delta$
\bewijsregel{\bigl|\bigl(f(x) + g(x)\bigr) - (L + M)\bigr| \leqslant |f(x) - L| + |g(x) - M| < \tfrac{\varepsilon}{2} + \tfrac{\varepsilon}{2} = \varepsilon .}
Voor het verschil schrijf je $-(g(x) - M)$ in plaats van $g(x) - M$; de absolute waarde blijft dezelfde.
\end{bewijs}
\end{theorie}

De limiet van een som is de som van de limieten. Let op de voorwaarde: beide limieten moeten bestaan.

\voorbeeld{}{Van twee functies weet je dat $\lim_{x \to 2} f(x) = 4$ en $\lim_{x \to 2} g(x) = -1$. Bereken $\lim_{x \to 2} \bigl(f(x) - g(x) + x\bigr)$.\label{vb:h1-somregel}}{%
\uitleg{De voorschriften kennen we niet, en dat hoeft ook niet: de somregel werkt alleen met de limieten. Elk van de drie termen heeft een limiet, dus we mogen splitsen.}
\begin{stappen}
\lim_{x \to 2} f(x) - \lim_{x \to 2} g(x) + \lim_{x \to 2} x & Somregel, twee keer toegepast: de limiet van de som is de som van de limieten. \\[1ex]
= 4 - (-1) + 2 = 7 & De eerste twee limieten zijn gegeven, de derde is een basislimiet.
\end{stappen}}

De voorwaarde doet ertoe. Neem $f(x) = \frac1x$ en $g(x) = -\frac1x$ in $a = 0$. Hun som is $0$, met limiet $0$, maar $\lim_{x \to 0} \frac1x$ bestaat niet. Je kunt die limiet dus niet splitsen in twee limieten.

\begin{oefeningen}
\oef{h1-somregel}
\end{oefeningen}

\subsection{Limiet van een product}

Bij een product verwacht je hetzelfde: $f(x)g(x)$ komt dicht bij $LM$. Het bewijs is wat langer, omdat een fout op $f(x)$ vermenigvuldigd wordt met $g(x)$.

\begin{theorie}
\begin{stelling}{Productregel}
Als $\lim_{x \to a} f(x) = L$ en $\lim_{x \to a} g(x) = M$, dan is \quad $\formule{\lim_{x \to a} f(x)\,g(x) = L\,M}$.
\end{stelling}
\begin{bewijs}
Neem $\varepsilon > 0$. Tel $L\,g(x)$ op en trek het af:
\bewijsregel{f(x)g(x) - LM = \bigl(f(x) - L\bigr)g(x) + L\bigl(g(x) - M\bigr).}
Kies $\delta$ zo klein dat voor $0 < |x - a| < \delta$ geldt:
\bewijsregel{|g(x) - M| < 1, \qquad |f(x) - L| < \frac{\varepsilon}{2(|M| + 1)}, \qquad |g(x) - M| < \frac{\varepsilon}{2(|L| + 1)} .}
Uit de eerste voorwaarde volgt $|g(x)| \leqslant |g(x) - M| + |M| < |M| + 1$. Met de driehoeksongelijkheid is dan
\bewijsregel{|f(x)g(x) - LM| < \frac{\varepsilon}{2(|M| + 1)} \cdot (|M| + 1) + |L| \cdot \frac{\varepsilon}{2(|L| + 1)} < \frac{\varepsilon}{2} + \frac{\varepsilon}{2} = \varepsilon .}
\end{bewijs}
\end{theorie}

Twee gevallen komen zo vaak voor dat je ze apart onthoudt. Is $g$ de constante $c$, dan geeft de productregel met de basislimiet $\lim_{x \to a} c\,f(x) = c\,L$: een constante factor mag voor de limiet. En omdat $x^n = x \cdot x \cdots x$, geeft de productregel $n - 1$ keer toegepast $\lim_{x \to a} x^n = a^n$.

\voorbeeld{}{Met dezelfde $f$ als in het vorige voorbeeld: bereken $\lim_{x \to 2} x^2 f(x)$.\label{vb:h1-productregel}}{%
\uitleg{De functie is een product van drie factoren: $x \cdot x \cdot f(x)$. Elke factor heeft een limiet, dus de productregel mag.}
\begin{stappen}
\lim_{x \to 2} x \cdot \lim_{x \to 2} x \cdot \lim_{x \to 2} f(x) & Productregel, twee keer. \\[1ex]
= 2 \cdot 2 \cdot 4 = 16 & Basislimieten en de gegeven limiet.
\end{stappen}}

\begin{oefeningen}
\oef{h1-productregel}
\end{oefeningen}

\subsection{Limiet van een quotiënt}

Bij een quotiënt komt er een voorwaarde bij: je mag niet door nul delen. Daarom moet de limiet van de noemer verschillend zijn van nul.

\begin{theorie}
\begin{stelling}{Quotiëntregel}
Als $\lim_{x \to a} f(x) = L$ en $\lim_{x \to a} g(x) = M$ met $M \neq 0$, dan is \quad $\formule{\lim_{x \to a} \frac{f(x)}{g(x)} = \frac{L}{M}}$.
\end{stelling}
\begin{bewijs}
We tonen eerst aan dat $\frac{1}{g(x)} \to \frac{1}{M}$. Neem $\varepsilon > 0$, en kies $\delta$ zo klein dat voor $0 < |x - a| < \delta$
\bewijsregel{|g(x) - M| < \frac{|M|}{2} \qquad\text{en}\qquad |g(x) - M| < \frac{\varepsilon M^2}{2} .}
Uit $|M| \leqslant |M - g(x)| + |g(x)| < \frac{|M|}{2} + |g(x)|$ volgt $|g(x)| > \frac{|M|}{2}$, dus $g(x) \neq 0$. Dan is
\bewijsregel{\Bigl|\frac{1}{g(x)} - \frac{1}{M}\Bigr| = \frac{|M - g(x)|}{|g(x)|\,|M|} < \frac{\varepsilon M^2 / 2}{\frac{|M|}{2} \cdot |M|} = \varepsilon .}
Nu is $\frac{f(x)}{g(x)} = f(x) \cdot \frac{1}{g(x)}$, en de productregel geeft de limiet $L \cdot \frac1M = \frac{L}{M}$.
\end{bewijs}
\end{theorie}

\voorbeeld{}{Met $\lim_{x \to 2} f(x) = 4$ en $\lim_{x \to 2} g(x) = -1$: bereken $\lim_{x \to 2} \dfrac{f(x) + 3}{g(x) + x}$.\label{vb:h1-quotientregel}}{%
\uitleg{Eerst controleren we de voorwaarde: de limiet van de noemer mag niet nul zijn.}
\begin{stappen}
\lim_{x \to 2} \bigl(g(x) + x\bigr) = -1 + 2 = 1 & Somregel in de noemer. Dat is niet nul, dus de quotiëntregel mag. \\[1ex]
\lim_{x \to 2} \bigl(f(x) + 3\bigr) = 4 + 3 = 7 & Somregel in de teller, met de basislimiet van de constante $3$. \\[1ex]
\lim_{x \to 2} \frac{f(x) + 3}{g(x) + x} = \frac71 = 7 & Quotiëntregel.
\end{stappen}}

Gaat de noemer naar $0$, dan zegt de quotiëntregel niets. Bij $\lim_{x \to 2} \frac{f(x)}{g(x) + 1}$ heb je te weinig gegevens: daar kan van alles uitkomen. Dat geval bekijken we in de volgende delen.

\begin{oefeningen}
\oef{h1-quotientregel}
\end{oefeningen}

\subsection{Invullen}

Uit de rekenregels volgt nu iets heel handigs: bij de functies die je het meest gebruikt, is de limiet gewoon de functiewaarde.

\begin{theorie}
\begin{stelling}{Limiet door invullen}
Is $p$ een veelterm, dan is \quad $\formule{\lim_{x \to a} p(x) = p(a)}$. \\
Is $\frac{p}{q}$ een rationale functie met $q(a) \neq 0$, dan is \quad $\formule{\lim_{x \to a} \frac{p(x)}{q(x)} = \frac{p(a)}{q(a)}}$.
\end{stelling}
\begin{bewijs}
Een veelterm is een som van termen $c\,x^k$. Uit de basislimieten en de productregel volgt $\lim_{x \to a} c\,x^k = c\,a^k$, en de somregel telt die limieten op tot $p(a)$. Voor een rationale functie is $\lim_{x \to a} q(x) = q(a) \neq 0$, en de quotiëntregel geeft $\frac{p(a)}{q(a)}$.
\end{bewijs}
Ook voor wortels, $\sin$, $\cos$, $\e^x$ en $\ln$ is de limiet in een punt van het domein de functiewaarde: je mag \textbf{\textcolor{accent}{invullen}}.
\end{theorie}

Het laatste bewijzen we niet: dat vraagt meer kennis van die functies dan we hier opbouwen. We gebruiken het als een gegeven.

\voorbeeld{}{Bereken $\displaystyle\lim_{x \to -1} \dfrac{x^3 + 2x}{x^2 + 1}$.\label{vb:h1-invullen}}{%
\uitleg{Dit is een rationale functie. Ligt $-1$ in het domein, dan mogen we invullen.}
\begin{stappen}
(-1)^2 + 1 = 2 \neq 0 & De noemer is niet nul in $-1$: het punt ligt in het domein. \\[1ex]
\lim_{x \to -1} \frac{x^3 + 2x}{x^2 + 1} = \frac{-1 - 2}{2} = -\frac32 & Invullen in teller en noemer.
\end{stappen}}

\begin{oefeningen}
\oef{h1-invullen}
\end{oefeningen}

\subsection{De onbepaalde vorm \texorpdfstring{$\frac00$}{0/0}}

Bij $\lim_{x \to 1} \frac{x^2 - 1}{x - 1}$ uit het begin van dit hoofdstuk lukt invullen niet: $1$ ligt niet in het domein, en de noemer gaat naar $0$. Ook de teller gaat naar $0$. Wat zegt dat over de limiet? Vergelijk drie functies waarbij teller en noemer allebei naar $0$ gaan:
\[\lim_{x \to 0} \frac{x}{x} = 1 , \qquad \lim_{x \to 0} \frac{x^2}{x} = \lim_{x \to 0} x = 0 , \qquad \lim_{x \to 0} \frac{x}{x^2} = \lim_{x \to 0} \frac1x \text{ bestaat niet.}\]
Uit \textit{teller en noemer gaan naar $0$} kun je dus niets besluiten.

\begin{theorie}
Geeft invullen $\frac00$, dan heet dat een \textbf{\textcolor{accent}{onbepaalde vorm}}: de limiet kan elk getal zijn, of niet bestaan.

\begin{stelling}{Gelijk buiten \texorpdfstring{$a$}{a}, gelijke limiet}
Is $f(x) = g(x)$ voor alle $x \neq a$ in een interval rond $a$, en is $\lim_{x \to a} g(x) = L$, dan is ook $\formule{\lim_{x \to a} f(x) = L}$.
\end{stelling}
\begin{bewijs}
Neem $\varepsilon > 0$ en de $\delta$ die bij $g$ hoort. Maak $\delta$ zo nodig kleiner, zodat $0 < |x - a| < \delta$ binnen het interval valt. Voor die $x$ is $f(x) = g(x)$, dus
\bewijsregel{|f(x) - L| = |g(x) - L| < \varepsilon .}
\end{bewijs}
\end{theorie}

Deze stelling verklaart waarom vereenvoudigen mag. Schrap je een factor $x - a$ in teller en noemer, dan verandert de functie alleen in $a$ zelf, en daar kijkt de limiet niet naar.

\voorbeeld{}{Bereken $\displaystyle\lim_{x \to 3} \dfrac{x^2 - 9}{x - 3}$.\label{vb:h1-ontbinden}}{%
\uitleg{Invullen geeft $\frac{9 - 9}{3 - 3} = \frac00$: een onbepaalde vorm. Teller en noemer zijn nul in $3$, dus ze bevatten allebei de factor $x - 3$.}
\begin{stappen}
\frac{x^2 - 9}{x - 3} = \frac{(x - 3)(x + 3)}{x - 3} & Ontbind de teller als verschil van twee kwadraten. \\[1ex]
= x + 3 \quad (x \neq 3) & Schrap de factor $x - 3$. Dat mag, want de limiet kijkt alleen naar $x \neq 3$. \\[1ex]
\lim_{x \to 3} (x + 3) = 6 & Nu is invullen wel mogelijk.
\end{stappen}}

\begin{oefeningen}
\oef{h1-ontbinden}
\end{oefeningen}

Ontbinden helpt niet altijd. Staat er een wortel, dan werk je die weg met het merkwaardig product
\[(a - b)(a + b) = a^2 - b^2 .\]

\voorbeeld{}{Bereken $\displaystyle\lim_{x \to 0} \dfrac{\sqrt{x + 4} - 2}{x}$.\label{vb:h1-wortel}}{%
\uitleg{Invullen geeft $\frac{2 - 2}{0} = \frac00$. De teller valt niet te ontbinden. Vermenigvuldig teller en noemer met de toegevoegde $\sqrt{x + 4} + 2$: dan verdwijnt de wortel uit de teller.}
\begin{stappenbreed}
\frac{(\sqrt{x + 4} - 2)(\sqrt{x + 4} + 2)}{x(\sqrt{x + 4} + 2)} = \frac{(x + 4) - 4}{x(\sqrt{x + 4} + 2)} & In de teller staat nu $a^2 - b^2$ met $a = \sqrt{x + 4}$ en $b = 2$. \\[1ex]
= \frac{x}{x(\sqrt{x + 4} + 2)} = \frac{1}{\sqrt{x + 4} + 2} & De factor $x$ verschijnt in de teller, en die schrappen we. \\[1ex]
\lim_{x \to 0} \frac{1}{\sqrt{x + 4} + 2} = \frac{1}{2 + 2} = \frac14 & Invullen: de noemer is nu $4$.
\end{stappenbreed}}

\begin{oefeningen}
\oef{h1-wortel}
\end{oefeningen}

\begin{werkwijze}
Vul in. Krijg je een getal (met een noemer verschillend van nul), dan is dat de limiet. Krijg je $\frac00$: ontbind teller en noemer en schrap de gemeenschappelijke factor, of vermenigvuldig met de toegevoegde bij een wortel. Vul dan opnieuw in.
\end{werkwijze}

Niet elke nul in de noemer is een onbepaalde vorm. Geeft invullen $\frac{3}{0}$, dan valt er niets te schrappen: de teller is niet nul. Zo'n limiet is geen getal. Wat er dan gebeurt, zie je in de volgende sectie.

\begin{oefeningen}
\oef{h1-berekenen}
\end{oefeningen}

%========================================================================================
\section{Oneindige limieten en asymptoten}

\subsection{Oneindige limieten}

\naast{%
Bij $f(x) = \frac{1}{x - 1}$ dicht bij $1$ gaat de teller naar $1$ en de noemer naar $0$. Het quotiënt wordt onbeperkt groot: rechts van $1$ positief, links negatief. Bij $x = 1{,}001$ is $f(x) = 1000$, bij $x = 0{,}999$ is $f(x) = -1000$.

Zo'n limiet is geen getal. Toch willen we het gedrag kort beschrijven, en daarvoor gebruiken we het symbool $\infty$.}{%
\begin{tikzpicture}[x=.55cm,y=.4cm,font=\scriptsize]
  \rooster{-2}{1}{4}{-4}{1}{4}
  \draw[grijs,dashed] (1,-4.3) -- (1,4.3);
  \kromme{-2}{.77}{1/(\x-1)}
  \kromme{1.23}{4}{1/(\x-1)}
\end{tikzpicture}}

\begin{theorie}
$\formule{\lim_{x \to a} f(x) = +\infty}$ betekent: $f(x)$ wordt groter dan elk getal als $x$ voldoende dicht bij $a$ komt. Nauwkeurig: voor elke $K > 0$ bestaat er een $\delta > 0$ zodat
\[\formule{0 < |x - a| < \delta \Rightarrow f(x) > K}.\]
Voor $-\infty$ eis je $f(x) < -K$. Eenzijdige oneindige limieten gebruiken alleen $x < a$ of alleen $x > a$.

\begin{stelling}{Getal gedeeld door nul}
Als $\lim_{x \to a} f(x) = L$ met $L \neq 0$, en $\lim_{x \to a} g(x) = 0$ met $g(x)$ steeds hetzelfde teken voor $x$ dicht bij $a$, dan is $\formule{\lim_{x \to a} \frac{f(x)}{g(x)} = +\infty}$ als $L$ en $g(x)$ hetzelfde teken hebben, en $\formule{-\infty}$ als ze een verschillend teken hebben.
\end{stelling}
\begin{bewijs}
Stel eerst $L > 0$ en $g(x) > 0$. Neem $K > 0$, en kies $\delta$ zo klein dat voor $0 < |x - a| < \delta$
\bewijsregel{|f(x) - L| < \frac{L}{2} \qquad\text{en}\qquad 0 < g(x) < \frac{L}{2K} .}
Dan is $f(x) > \frac{L}{2}$, en
\bewijsregel{\frac{f(x)}{g(x)} > \frac{L/2}{L/(2K)} = K .}
Dus $\frac{f(x)}{g(x)} \to +\infty$. De andere gevallen volgen hieruit. Is $L < 0$, pas dit dan toe op $-f$: dan gaat $-\frac{f(x)}{g(x)}$ naar $+\infty$. Is $g(x) < 0$, schrijf dan $\frac{f(x)}{g(x)} = \frac{-f(x)}{-g(x)}$. Voor eenzijdige limieten werkt het bewijs net zo, met alleen $x < a$ of $x > a$.
\end{bewijs}
\end{theorie}

In het kort: \textit{getal gedeeld door nul} geeft $\pm\infty$, en de tekens van teller en noemer bepalen welk. Wisselt de noemer in $a$ van teken, dan zijn de linker- en de rechterlimiet verschillend.

\voorbeeld{}{Bereken $\liml{x}{2} \dfrac{x + 1}{x - 2}$ en $\limr{x}{2} \dfrac{x + 1}{x - 2}$.\label{vb:h1-oneindig}}{%
\naast{%
\uitleg{Invullen geeft $\frac30$: geen onbepaalde vorm, maar een getal gedeeld door nul. Het teken van de noemer bepaalt het teken van het resultaat, en dat verschilt links en rechts van $2$.}
\begin{stappenvrij}{3.4cm}{6.2cm}
x + 1 \to 3 > 0 & De teller is positief. \\
\liml{x}{2} \frac{x + 1}{x - 2} = -\infty & Voor $x < 2$ is $x - 2 < 0$: positief gedeeld door negatief. \\[1.5ex]
\limr{x}{2} \frac{x + 1}{x - 2} = +\infty & Voor $x > 2$ is $x - 2 > 0$.
\end{stappenvrij}}{%
\begin{tikzpicture}[x=.45cm,y=.3cm,font=\scriptsize]
  \rooster{-2}{1}{6}{-6}{2}{8}
  \draw[grijs,dashed] (2,-6.3) -- (2,8.3);
  \kromme{-2}{1.67}{(\x+1)/(\x-2)}
  \kromme{2.4}{6}{(\x+1)/(\x-2)}
\end{tikzpicture}}}

\begin{oefeningen}
\oef{h1-eenzijdig}
\end{oefeningen}

\subsection{Verticale asymptoten}

In de figuur bij het vorige voorbeeld klimt de grafiek tegen de rechte $x = 2$ op, zonder ze ooit te raken. Zo'n rechte heet een asymptoot.

\begin{theorie}
De rechte $x = a$ is een \textbf{\textcolor{accent}{verticale asymptoot}} van de grafiek van $f$ als minstens één eenzijdige limiet in $a$ gelijk is aan $+\infty$ of $-\infty$.
\end{theorie}

Een nulwaarde van de noemer geeft niet altijd een asymptoot. Is de teller er ook nul, dan heb je een onbepaalde vorm, en die kan een gewoon getal opleveren: dan heeft de grafiek een gat, een \textit{perforatie}.

\voorbeeld{}{Bepaal de verticale asymptoten van de grafiek van $f(x) = \dfrac{x^2 - x}{x^2 - 1}$.\label{vb:h1-va}}{%
\naast{%
\uitleg{De kandidaten zijn de nulwaarden van de noemer, $1$ en $-1$. Kijk telkens wat invullen in de teller geeft.}
\begin{stappenvrij}{4cm}{5.6cm}
f(x) = \frac{x(x - 1)}{(x - 1)(x + 1)} & Ontbind teller en noemer. \\
{\frac{-1 - 1}{0} = \frac{-2}{0}} & In $-1$: een getal gedeeld door nul. De rechte $x = -1$ is een verticale asymptoot. \\
{\frac{1 - 1}{0} = \frac00} & In $1$: onbepaald. Schrap $x - 1$. \\
{\lim_{x \to 1} \frac{x}{x + 1} = \frac12} & Een getal: in $\bigl(1, \frac12\bigr)$ heeft de grafiek een gat, geen asymptoot.
\end{stappenvrij}}{%
\begin{tikzpicture}[x=.45cm,y=.35cm,font=\scriptsize]
  \rooster{-5}{1}{3}{-3}{1}{5}
  \draw[grijs,dashed] (-1,-3.3) -- (-1,5.3);
  \kromme{-5}{-1.5}{\x/(\x+1)}
  \kromme{-.75}{3}{\x/(\x+1)}
  \open{(1,.5)}
\end{tikzpicture}}}

\begin{werkwijze}
Zoek de nulwaarden van de noemer. Vul ze in de teller in. Is de teller niet nul, dan is er een verticale asymptoot; de eenzijdige limieten tonen hoe de grafiek ernaast loopt. Is de teller ook nul, schrap dan de gemeenschappelijke factor en vul opnieuw in.
\end{werkwijze}

\begin{oefeningen}
\oef{h1-va}
\end{oefeningen}

\subsection{Limieten in oneindig}

Ook als $x$ onbeperkt groot wordt, kunnen de functiewaarden een vast getal naderen. Bij $\frac1x$ zie je dat meteen: $\frac{1}{1000} = 0{,}001$ en $\frac{1}{10^6} = 0{,}000\,001$.

\begin{theorie}
$\formule{\lim_{x \to +\infty} f(x) = L}$ betekent: $f(x)$ komt willekeurig dicht bij $L$ als $x$ voldoende groot wordt. Nauwkeurig: voor elke $\varepsilon > 0$ bestaat er een getal $N$ zodat
\[\formule{x > N \Rightarrow |f(x) - L| < \varepsilon}.\]
Voor $x \to -\infty$ eis je $x < -N$.

\begin{stelling}{Basislimiet in oneindig}
Voor elk natuurlijk getal $n \geqslant 1$ is \quad $\formule{\lim_{x \to +\infty} \frac{1}{x^n} = 0}$ \quad en \quad $\formule{\lim_{x \to -\infty} \frac{1}{x^n} = 0}$.
\end{stelling}
\begin{bewijs}
Neem $\varepsilon > 0$ en $N$ het grootste van de getallen $1$ en $\frac{1}{\varepsilon}$. Voor $|x| > N$ is $|x| > 1$, dus $|x|^n \geqslant |x|$, en
\bewijsregel{\Bigl|\frac{1}{x^n} - 0\Bigr| = \frac{1}{|x|^n} \leqslant \frac{1}{|x|} < \frac{1}{N} \leqslant \varepsilon .}
Dat geldt zowel voor $x > N$ als voor $x < -N$.
\end{bewijs}
\end{theorie}

De rekenregels gelden ook voor $x \to \pm\infty$. De bewijzen zijn dezelfde, met een grens $N$ in plaats van een afstand $\delta$.

\voorbeeld{}{Bereken $\displaystyle\lim_{x \to +\infty} \dfrac{3x^2 - x}{2x^2 + 5}$.\label{vb:h1-oneindig2}}{%
\uitleg{Teller en noemer worden allebei onbeperkt groot: $\frac{\infty}{\infty}$ is ook een onbepaalde vorm. Deel teller en noemer door de hoogste macht van $x$ in de noemer: dan ontstaan termen die naar $0$ gaan.}
\begin{stappenbreed}
\frac{3x^2 - x}{2x^2 + 5} = \frac{3 - \frac1x}{2 + \frac{5}{x^2}} & Teller en noemer gedeeld door $x^2$. Het quotiënt verandert niet. \\[1ex]
\lim_{x \to +\infty} \frac{3 - \frac1x}{2 + \frac{5}{x^2}} = \frac{3 - 0}{2 + 0} = \frac32 & Basislimieten in oneindig, en dan de som- en de quotiëntregel.
\end{stappenbreed}}

\begin{oefeningen}
\oef{h1-oneindig}
\end{oefeningen}

Het resultaat hangt af van de graad van teller en noemer. Het volgende voorbeeld zet de drie mogelijkheden naast elkaar.

\voorbeeld{}{Bereken $\displaystyle\lim_{x \to +\infty}$ van $\dfrac{2x + 1}{x^2 + 3}$, van $\dfrac{2x^2 + 1}{x^2 + 3}$ en van $\dfrac{2x^3 + 1}{x^2 + 3}$.\label{vb:h1-graad}}{%
\uitleg{Telkens delen we door $x^2$, de hoogste macht in de noemer. De noemer wordt dan $1 + \frac{3}{x^2}$, met limiet $1$. Alleen de teller verschilt.}
\begin{stappenbreed}
\frac{\frac2x + \frac{1}{x^2}}{1 + \frac{3}{x^2}} \to \frac{0}{1} = 0 & Graad teller kleiner dan graad noemer: de limiet is $0$. \\[1.5ex]
\frac{2 + \frac{1}{x^2}}{1 + \frac{3}{x^2}} \to \frac21 = 2 & Gelijke graad: het quotiënt van de hoogstegraadscoëfficiënten. \\[1.5ex]
\frac{2x + \frac{1}{x^2}}{1 + \frac{3}{x^2}} \to +\infty & Graad teller groter: de teller wordt onbeperkt groot, de noemer gaat naar $1$.
\end{stappenbreed}}

\begin{werkwijze}
Voor een rationale functie in $\pm\infty$: deel teller en noemer door de hoogste macht van $x$ in de noemer, en gebruik $\frac{1}{x^n} \to 0$. Is de graad van de teller groter, bepaal dan het teken van $\pm\infty$ uit de hoogstegraadstermen.
\end{werkwijze}

\begin{oefeningen}
\oef{h1-graad}
\end{oefeningen}

\subsection{Horizontale asymptoten}

Nadert $f(x)$ een getal $L$ als $x$ onbeperkt groot wordt, dan komt de grafiek steeds dichter bij de horizontale rechte $y = L$.

\begin{theorie}
De rechte $y = L$ is een \textbf{\textcolor{accent}{horizontale asymptoot}} van de grafiek van $f$ als $\formule{\lim_{x \to +\infty} f(x) = L}$ of $\formule{\lim_{x \to -\infty} f(x) = L}$.
\end{theorie}

\voorbeeld{}{Bepaal met limieten de asymptoten van de grafiek van $f(x) = \dfrac{2x + 1}{x - 1}$.\label{vb:h1-ha}}{%
\naast{%
\uitleg{Een verticale asymptoot zoek je bij de nulwaarden van de noemer, een horizontale met de limieten in $\pm\infty$.}
\begin{stappenvrij}{4cm}{5.6cm}
{\frac{2 + 1}{1 - 1} = \frac30} & In $1$: getal gedeeld door nul, dus verticale asymptoot $x = 1$. \\
\liml{x}{1} f(x) = -\infty & Teller positief, noemer negatief. \\
\limr{x}{1} f(x) = +\infty & Noemer positief. \\
\lim_{x \to \pm\infty} \frac{2 + \frac1x}{1 - \frac1x} = 2 & Gelijke graad: horizontale asymptoot $y = 2$, aan beide kanten.
\end{stappenvrij}}{%
\begin{tikzpicture}[x=.45cm,y=.35cm,font=\scriptsize]
  \rooster{-4}{1}{5}{-3}{1}{7}
  \draw[grijs,dashed] (1,-3.3) -- (1,7.3) (-4,2) -- (5,2);
  \kromme{-4}{.4}{2+3/(\x-1)}
  \kromme{1.6}{5}{2+3/(\x-1)}
\end{tikzpicture}}}

\begin{oefeningen}
\oef{h1-asymptoten}
\end{oefeningen}

%========================================================================================
\section{Continuïteit}

\subsection{Continu in een punt}

Een grafiek die je kunt tekenen zonder je pen op te heffen, maakt geen sprongen en heeft geen gaten. Waar dat misloopt, zie je in de figuren: er is een gat, een sprong, of een punt dat niet op zijn plaats ligt.

\figuren{\begin{tikzpicture}[x=.5cm,y=.5cm,font=\scriptsize]
  \roosterleeg{-1}{1}{3}{-1}{1}{3}
  \kromme{-1}{3}{.25*(\x)^2+1}
  \open{(1,1.25)}
  \node[below,inner sep=1.5pt] at (1,0) {$a$};
  \node[above] at (2.2,3.4) {gat};
\end{tikzpicture}\qquad
\begin{tikzpicture}[x=.5cm,y=.5cm,font=\scriptsize]
  \roosterleeg{-1}{1}{3}{-1}{1}{3}
  \kromme{-1}{1}{.5*\x+.5}
  \kromme{1}{3}{.5*\x+1.5}
  \vol{(1,1)} \open{(1,2)}
  \node[below,inner sep=1.5pt] at (1,0) {$a$};
  \node[above] at (2.2,3.4) {sprong};
\end{tikzpicture}\qquad
\begin{tikzpicture}[x=.5cm,y=.5cm,font=\scriptsize]
  \roosterleeg{-1}{1}{3}{-1}{1}{3}
  \kromme{-1}{3}{.25*(\x)^2+1}
  \open{(1,1.25)} \vol{(1,2.5)}
  \node[below,inner sep=1.5pt] at (1,0) {$a$};
  \node[above] at (2,3.4) {punt ernaast};
\end{tikzpicture}}

In de eerste figuur bestaat $f(a)$ niet. In de tweede bestaat de limiet niet. In de derde bestaan ze allebei, maar ze zijn niet gelijk. De definitie sluit alle drie uit.

\begin{theorie}
Een functie $f$ is \textbf{\textcolor{accent}{continu}} in $a$ als $\formule{\lim_{x \to a} f(x) = f(a)}$: de limiet bestaat, $f(a)$ bestaat, en ze zijn gelijk. Ze is continu op een interval als ze continu is in elk punt ervan.
\end{theorie}

\voorbeeld{}{Neem $f(x) = \dfrac{x^2 - 1}{x - 1}$ voor $x \neq 1$ en $f(1) = 2$. Is $f$ continu in $1$?\label{vb:h1-continupunt}}{%
\uitleg{Controleer de drie voorwaarden één voor één.}
\begin{stappen}
f(1) = 2 & De functiewaarde bestaat: ze is apart gegeven. \\[1ex]
\lim_{x \to 1} \frac{(x - 1)(x + 1)}{x - 1} = \lim_{x \to 1} (x + 1) = 2 & De limiet bestaat: onbepaalde vorm, ontbinden en schrappen. \\[1ex]
\lim_{x \to 1} f(x) = f(1) & Ze zijn gelijk: $f$ is continu in $1$.
\end{stappen}
Dit is de functie uit het begin van het hoofdstuk, met het gat opgevuld. Kies je $f(1) = 3$, dan is ze niet continu: dan krijg je de derde figuur.}

\begin{oefeningen}
\oef{h1-continupunt}
\end{oefeningen}

\subsection{Continue functies}

Moet je continuïteit telkens met de definitie nagaan? Nee: uit de rekenregels voor limieten volgen meteen rekenregels voor continue functies.

\begin{theorie}
\begin{stelling}{Bewerkingen met continue functies}
Zijn $f$ en $g$ continu in $a$, dan zijn ook $\formule{f + g}$, $\formule{f - g}$ en $\formule{f\,g}$ continu in $a$, en $\formule{\frac{f}{g}}$ als $g(a) \neq 0$.
\end{stelling}
\begin{bewijs}
Gegeven is $\lim_{x \to a} f(x) = f(a)$ en $\lim_{x \to a} g(x) = g(a)$. De somregel geeft
\bewijsregel{\lim_{x \to a} \bigl(f(x) + g(x)\bigr) = f(a) + g(a) ,}
en dat is de waarde van $f + g$ in $a$. Voor het verschil, het product en het quotiënt gebruik je de andere rekenregels; bij het quotiënt is de voorwaarde $g(a) \neq 0$ precies die van de quotiëntregel.
\end{bewijs}
\end{theorie}

Daarom zijn veeltermen overal continu, en rationale functies in elk punt van hun domein: dat zegt de stelling over invullen. Ook $\sin$, $\cos$ en $\e^x$ zijn overal continu, en wortels en $\ln$ op hun domein. Problemen kunnen er dus alleen zijn waar een functie niet bestaat, of waar ze met verschillende voorschriften is gedefinieerd: op de overgang kan ze een sprong maken.

\voorbeeld{}{Bepaal $k$ zodat $f$ continu is: $f(x) = x^2$ voor $x \leqslant 1$ en $f(x) = 2x + k$ voor $x > 1$.\label{vb:h1-continu}}{%
\naast{%
\uitleg{Links en rechts van $1$ is $f$ een veelterm, dus continu. Alleen in $1$ kan het misgaan: daar moeten de linker- en de rechterlimiet gelijk zijn aan $f(1)$.}
\begin{stappenvrij}{5.3cm}{4.4cm}
f(1) = 1^2 = 1 & Voor $x = 1$ geldt het eerste voorschrift. \\
{\liml{x}{1} f(x) = \liml{x}{1} x^2 = 1} & Links van $1$: het eerste voorschrift. \\
{\limr{x}{1} f(x) = \limr{x}{1} (2x + k) = 2 + k} & Rechts: het tweede. \\
2 + k = 1 \iff k = -1 & Gelijkstellen: geen sprong meer.
\end{stappenvrij}}{%
\begin{tikzpicture}[x=.55cm,y=.45cm,font=\scriptsize]
  \rooster{-2}{1}{3}{-1}{1}{5}
  \kromme{-2}{1}{(\x)^2}
  \kromme{1}{3}{2*\x-1}
  \vol{(1,1)}
\end{tikzpicture}}}

\begin{werkwijze}
Bij een functie met meerdere voorschriften: controleer elk voorschrift apart, en dan elk overgangspunt $a$. Bereken daar $f(a)$, de linkerlimiet en de rechterlimiet met het juiste voorschrift, en stel ze gelijk.
\end{werkwijze}

\begin{oefeningen}
\oef{h1-continuk}
\end{oefeningen}

\subsection{De tussenwaardestelling}

Een continue functie die onder de $x$-as begint en erboven eindigt, moet de $x$-as ergens snijden. Dat lijkt vanzelfsprekend, en het is een van de belangrijkste eigenschappen van continue functies.

\needspace{12\baselineskip}
\begin{theorie}
\begin{stelling}{Tussenwaardestelling}
\begin{minipage}[t]{.6\linewidth}
\vspace{0pt}
Is $f$ continu op $[a, b]$ en hebben $f(a)$ en $f(b)$ een verschillend teken, dan heeft $f$ een nulwaarde in $\left]a, b\right[$: \\[1mm]
\qquad $\formule{\text{er is een } c \in \left]a, b\right[ \text{ met } f(c) = 0}$.
\end{minipage}\hfill
\begin{minipage}[t]{.34\linewidth}
\vspace{0pt}
\begin{tikzpicture}[x=.9cm,y=.5cm,font=\scriptsize]
  \draw[->,grijs] (-.1,0) -- (3.4,0) node[right,text=black] {$x$};
  \kromme{.3}{3}{.5*(\x-1.6)^3+.3*(\x-1.6)+.2}
  \vol{(.3,-1.29)} \vol{(3,1.99)}
  \draw[grijs,dashed] (.3,0) -- (.3,-1.29) (3,0) -- (3,1.99);
  \node[above] at (.3,0) {$a$};
  \node[below] at (3,0) {$b$};
  \fill (1.12,0) circle (1.6pt) node[above left] {$c$};
\end{tikzpicture}
\end{minipage}
\end{stelling}
\end{theorie}

Het bewijs steunt op een eigenschap van de reële getallen: de getallenrechte heeft geen gaten. Die eigenschap behandelen we niet, dus we laten het bewijs weg.

De stelling zegt dat er minstens één nulwaarde is. Ze zegt niet hoeveel, en ook niet waar. De voorwaarde \textit{continu} is nodig: $f(x) = \frac1x$ op $[-1, 1]$ heeft $f(-1) = -1$ en $f(1) = 1$, maar geen nulwaarde. Daar is geen tegenspraak, want $\frac1x$ bestaat niet in $0$: de grafiek springt over de $x$-as.

Je kunt de stelling ook gebruiken voor andere waarden dan $0$. Ligt een getal $d$ tussen $f(a)$ en $f(b)$, pas haar dan toe op $f(x) - d$: er is een $c$ met $f(c) = d$. Een continue functie neemt dus elke tussenliggende waarde aan, vandaar de naam.

\voorbeeld{}{Toon aan dat $x^3 + x - 1 = 0$ een oplossing heeft tussen $0$ en $1$. Zoek ze op één decimaal.\label{vb:h1-tussenwaarde}}{%
\uitleg{Neem $f(x) = x^3 + x - 1$: een veelterm, dus continu. Zoek een tekenwissel, en halveer daarna het interval telkens, zodat de oplossing steeds nauwer ingesloten wordt.}
\begin{stappenbreed}
f(0) = -1 < 0,\quad f(1) = 1 > 0 & Tekenwissel: volgens de tussenwaardestelling ligt er een nulwaarde in $\left]0, 1\right[$. \\[1ex]
f(0{,}5) = -0{,}375 < 0 & Het teken wisselt nu tussen $0{,}5$ en $1$. \\[1ex]
f(0{,}75) \approx 0{,}17 > 0 & De nulwaarde ligt tussen $0{,}5$ en $0{,}75$. \\[1ex]
f(0{,}65) \approx -0{,}08 < 0,\quad f(0{,}7) = 0{,}043 > 0 & Ze ligt tussen $0{,}65$ en $0{,}7$: op één decimaal is ze $0{,}7$.
\end{stappenbreed}}

\begin{werkwijze}
Herschrijf de vergelijking als $f(x) = 0$ en controleer dat $f$ continu is. Zoek twee getallen waar $f$ een verschillend teken heeft. Bereken $f$ in een punt ertussen en houd de helft met de tekenwissel. Herhaal tot het interval klein genoeg is.
\end{werkwijze}

\begin{oefeningen}
\oef{h1-continu}

\oef{h1-tws}
\end{oefeningen}

%========================================================================================
\opdrachtenkop

\begin{oefeningen}
\oef{h1-o1}

\needspace{14\baselineskip}
\oef{h1-o10}

\needspace{16\baselineskip}
\oef{h1-o2}

\needspace{10\baselineskip}
\oef{h1-o14}

\needspace{10\baselineskip}
\oef{h1-o11}

\needspace{12\baselineskip}
\oef{h1-o15}

\needspace{16\baselineskip}
\oef{h1-o3}

\needspace{10\baselineskip}
\oef{h1-o12}

\needspace{10\baselineskip}
\oef{h1-o13}

\oef{h1-o4}

\needspace{16\baselineskip}
\oef{h1-o16}

\oef{h1-o5}

\needspace{12\baselineskip}
\oef{h1-o17}

\needspace{10\baselineskip}
\oef{h1-o18}

\needspace{10\baselineskip}
\oef{h1-o6}

\needspace{12\baselineskip}
\oef{h1-o19}

\oef{h1-o7}

\needspace{8\baselineskip}
\oef{h1-o21}

\needspace{10\baselineskip}
\oef{h1-o22}

\needspace{8\baselineskip}
\oef{h1-o20}

\needspace{8\baselineskip}
\oef{h1-o23}

\needspace{8\baselineskip}
\oef{h1-o24}

\oef{h1-o8}

\oef{h1-o9}

\needspace{10\baselineskip}
\oef{h1-o25}

\needspace{10\baselineskip}
\oef{h1-o28}

\needspace{14\baselineskip}
\oef{h1-o26}

\needspace{10\baselineskip}
\oef{h1-o27}
\end{oefeningen}

%----------------------------------------------------------------------------------------
% KOPPEL: verbanden met andere leerstof
%----------------------------------------------------------------------------------------
\needspace{20\baselineskip}
\opdrachtsectie{Koppel}

\lkr{Rationale functies, Rijen en Exponentiële en logaritmische functies: kijk na wat al gezien is. Het pakket legt de volgorde niet vast.}

\begin{oefeningen}
\oef{h1-koppel-1}

\oef{h1-koppel-2}

\oef{h1-koppel-3}
\end{oefeningen}

%----------------------------------------------------------------------------------------
\opdrachtsectie{Test jezelf}

\begin{oefeningen}
\oef{h1-t6}

\oef{h1-t7}

\oef{h1-t1}

\oef{h1-t2}

\oef{h1-t3}

\oef{h1-t4}

\oef{h1-t5}

\oef{h1-t8}
\end{oefeningen}

\herhaalschema{%
  \fv{oef:h1-t6}{vb:h1-aflezen}{oef:h1-o10}%
  \fv{oef:h1-t7}{vb:h1-definitie}{oef:h1-begrip}%
  \fv{oef:h1-t1}{vb:h1-ontbinden}{oef:h1-ontbinden}%
  \fv{oef:h1-t2}{vb:h1-wortel}{oef:h1-wortel}%
  \fv{oef:h1-t3}{vb:h1-oneindig}{oef:h1-eenzijdig}%
  \fv{oef:h1-t4}{vb:h1-oneindig2}{oef:h1-oneindig}%
  \fv{oef:h1-t5}{vb:h1-continu}{oef:h1-continuk}%
  \fv{oef:h1-t8}{vb:h1-tussenwaarde}{oef:h1-continu}}

\ifzelfstudie
\section*{Maak je eigen samenvatting}
Sluit dit boek en schrijf in je schrift: wat $\lim_{x \to a} f(x) = L$ betekent (in woorden en met $\varepsilon$ en $\delta$), linker- en rechterlimiet, de rekenregels, wat je doet bij $\frac00$ en bij $\frac{\infty}{\infty}$, de definities van verticale en horizontale asymptoot, wat continu betekent, en de tussenwaardestelling. Vergelijk daarna met de kaders.
\fi
